% from main.tex
%%%%%%%%%%  scan idx 268  |  folio 255  %%%%%%%%%%
preuve de (ii) \(\Longrightarrow\) (i).

(i) \(\Longrightarrow\) (ii). La question étant locale au voisinage de \(x\), on peut supposer que \(f\) se factorise en \(X\overset{i}{\hookrightarrow}X'\overset{f'}{\longrightarrow}Y\), où \(i\) est une immersion fermée et \(f'\) un morphisme lisse. Il est immédiat que les conditions (i) et (ii) ne sont pas altérées lorsqu'on remplace \(E\) par \(i_*(E)\), \(f\) par \(f'\) (et \(h\) par \(gf'\)), de sorte qu'on peut supposer \(f\) lisse. Comme \(g\) est parfait (4.3.1), \(h\) est parfait, donc (1.12) \(E\) est pseudo-cohérent. Quitte à se localiser, on peut, d'après (II 2.2.10 b)) et (I 2.10 a)), supposer que \(E\) est strictement pseudo-cohérent (I 2.1) et nul en degré \(>b\), et que la flèche de troncation
\begin{equation}
 E\longrightarrow(0\longrightarrow E^a/B^a\longrightarrow E^{a+1}\longrightarrow\cdots\longrightarrow E^b\longrightarrow0)\tag{$t$}
\end{equation}
est un quasi-isomorphisme (l'hypothèse (i) implique en effet que \(E\) est, au voisinage de \(s\), acyclique en dehors de \([a,b]\)). Il s'agit de prouver que \(E^a/B^a\) est \(f\)-plat au voisinage de \(x\). Notons \(i:X_s\longrightarrow X\) la flèche canonique. On a :
\begin{equation}
 E_x\otimes_{\mathcal O_{Y,y}}^{\mathbf L}\mathcal O_{Y_s,y}\xrightarrow{\sim}i^*(E)_x.\tag{$*$}
\end{equation}
Comme \(h\) est plat, l'hypothèse (i) implique, d'après (3.5.2), \(\tamp_{h_s}(i^*(E))\subset[a,b]\) au voisinage de \(x\). Donc quitte à se localiser, on peut supposer que la flèche de troncation
\begin{equation}
 i^*(E)\longrightarrow(0\longrightarrow i^*(E^a)/B^a(i^*(E))\longrightarrow\cdots\longrightarrow i^*(E^b)\longrightarrow0)\tag{$i^*(t)$}
\end{equation}
est un quasi-isomorphisme. L'hypothèse \(\tamp_{\mathcal O_{Y_s,y}}(i^*(E)_x)\subset[a,b]\) entraîne alors d'après (I 5.1.1), que \(i^*(E^a/B^a)_x\) est un \(\mathcal O_{Y_s,y}\)-module plat. D'autre part, la flèche \((*)\) étant un quasi-isomorphisme, l'hypothèse \(\tamp_h(E)\subset[a,b]\) implique, toujours d'après (I 5.1.1), que \(E^a/B^a\) est \(h\)-
